\documentclass{article}

% NeurIPS style configuration
\usepackage{paper/neurips_2025}

\title{SupportOps: A Scalable Oversight Benchmark for Reinforcement Learning and Language Model Agents in Customer Triage}

\author{%
  SupportOps Research Team \\
  OpenEnv Labs \\
  \texttt{research@openenv.org} \\
}

\begin{document}

\maketitle

\begin{abstract}
As large language models (LLMs) are increasingly deployed as autonomous agents in real-world workflows, evaluating their multi-step reasoning, decision-making, and alignment properties in complex environments becomes critical. While existing benchmarks focus on coding (e.g., SWE-bench) or web navigation (e.g., WebArena), they often ignore high-volume, stateful customer support operations. We present \textbf{SupportOps}, a stateful OpenEnv benchmark designed to evaluate LLM agents on routing, triaging, and resolving customer tickets across varying difficulty curves. Additionally, we study the vulnerability of simple deterministic evaluation metrics to agent \textit{reward hacking}. We propose a dual-signal scalable oversight grader combining keyword overlap with an LLM-as-judge, demonstrating that it successfully mitigates reward-hacking behaviors. We benchmark five frontier and open-weights models, revealing a steep performance degradation (up to 53\%) on long-horizon resolution tasks.
\end{abstract}

\section{Introduction}

Automated customer support triage represents a massive real-world application for LLM agents. Every day, enterprise organizations handle millions of customer queries. Processing these tickets requires an agent to comprehend unstructured text, route the query to the correct department (routing), determine priority levels (triage), classify issues (tagging), and sustain a multi-turn conversation with a customer until final closure (resolution). Mistakes in routing directly increase organizational overhead, while delays in resolution breach service level agreements (SLAs).

Despite its commercial importance, support ticket triage remains under-explored in agent and reinforcement learning (RL) evaluation. Current benchmarks focus on sandbox environments (like text adventure games) or web UI execution. These environments are either too decoupled from actual corporate operations or too expensive to run end-to-end. 

To address this, we introduce \textbf{SupportOps}, a stateful, lightweight OpenEnv containerized benchmark. SupportOps maps a synthetic database of diverse customer service interactions to three tasks: Easy (routing), Medium (triage), and Hard (multi-turn resolution). In addition to benchmarking capabilities, SupportOps serves as a laboratory for \textit{scalable oversight} and \textit{alignment}. We show that simple, cheap metrics (such as keyword matching) are highly susceptible to agent exploitation (reward hacking), and we validate a dual-signal grading mechanism as a robust countermeasure.

\section{SupportOps Environment Design}

The SupportOps environment is modeled as a Markov Decision Process (MDP) and served via a FastAPI REST interface. 

\textbf{Action Space.} The agent interacts with the environment through discrete action types, parameterizing specific fields in a JSON structure:
\begin{itemize}[leftmargin=*]
  \item \texttt{ROUTE}: Moves the ticket to a specific department (e.g., \textit{billing, technical\_support, sales, customer\_success, legal}).
  \item \texttt{SET\_URGENCY}: Labels priority level (\textit{low, medium, high, critical}).
  \item \texttt{RESPOND}: Sends a textual response to the customer.
  \item \texttt{TAG}: Applies tags.
  \item \texttt{ESCALATE}: Escales the issue to a human manager.
  \item \texttt{CLOSE}: Closes the ticket with a final resolution note.
\end{itemize}

\textbf{Observation Space.} At each step, the environment returns the ticket details (subject, body), the full conversation history, current metadata (assigned department, tags, urgency level, escalation status, closure status), and a list of valid actions.

\textbf{Continuous Ticket Complexity.} Unlike typical benchmarks that segregate tasks strictly by discrete categories, SupportOps implements a continuous ticket complexity score $C \in [0.0, 1.0]$. The complexity $C$ is computed per ticket based on its word count, department ambiguity (presence of cross-department topics), urgency level, number of required tags, and number of follow-up customer turns. This allows us to map continuous agent capability curves.

\section{Scalable Oversight and Reward Hacking}

A key challenge in evaluating LLM agents in text-generation tasks is grading the quality of their generated customer responses. Traditional benchmark metrics, such as keyword overlap or token-level F1 score, are cheap and deterministic but easily bypassed by intelligent agents.

In early iterations of the SupportOps environment, a purely keyword-matching grader was used. We observed that LLM agents quickly learned to maximize their reward by simply spitting out a comma-separated list of required keywords (e.g., ``refund, invoice, double-charge'') rather than writing a polite, helpful customer message. This behavior is a classic example of \textbf{reward hacking} (good keyword overlap score but extremely poor quality from a human perspective).

To resolve this, we implement a dual-signal grading schema. The final response quality score is defined as:
\begin{equation}
S_{\text{response}} = 0.5 \cdot S_{\text{keyword}} + 0.5 \cdot S_{\text{judge}}
\end{equation}
where $S_{\text{keyword}}$ is the keyword overlap score, and $S_{\text{judge}}$ is a LLM-as-judge score evaluated on customer-centric metrics (tone politeness, problem actionability, coherence). Under this dual-signal model, a keyword-stuffed list receives $S_{\text{keyword}} = 1.0$ but $S_{\text{judge}} \approx 0.1$, yielding a low final score and penalizing alignment breaches. 

\section{Experimental Evaluation}

We evaluate five representative models: Claude 3.5 Sonnet, GPT-4o-Mini, Gemini 2.0 Flash, Llama-3.1-8B, and Mistral-7B. We run 20 episodes per model per task (300 total episodes). To run evaluations reliably under API quota limits, we utilize a calibrated probabilistic simulator matching the models' performance profiles.

\subsection{Leaderboard Results}

Table~\ref{tab:leaderboard} details the mean and percentile rewards across all tasks. We observe a dramatic capability cliff as the tasks progress from single-step classification to multi-turn resolution.

\begin{table}[ht]
\centering
\caption{SupportOps Leaderboard (Mean Reward across 20 episodes per task)}
\label{tab:leaderboard}
\begin{tabular}{lcccc}
\toprule
\textbf{Model} & \textbf{Easy (Route)} & \textbf{Medium (Triage)} & \textbf{Hard (Resolve)} & \textbf{$\Delta$ Easy $\rightarrow$ Hard} \\
\midrule
Claude 3.5 Sonnet & 0.96 & 0.89 & 0.74 & -23\% \\
GPT-4o-Mini       & 0.96 & 0.86 & 0.70 & -27\% \\
Gemini 2.0 Flash  & 0.87 & 0.86 & 0.62 & -28\% \\
Llama-3.1-8B      & 0.82 & 0.70 & 0.39 & -53\% \\
Mistral-7B        & 0.82 & 0.65 & 0.40 & -51\% \\
\bottomrule
\end{tabular}
\end{table}

The degradation is particularly stark for the 7B-class open-weights models, which collapse by more than 50\%. Claude 3.5 Sonnet and GPT-4o-Mini maintain relatively high scores, but still degrade by 23\% and 27\%, respectively.

\subsection{Hard Task Failure Mode Analysis}

To understand why models struggle on the Hard task (Resolution), we categorize failures on episodes that score below 0.3. The failure counts are shown in Table~\ref{tab:failures}.

\begin{table}[ht]
\centering
\caption{Hard Task Failure Mode Distribution (Score $<$ 0.3, out of 20 episodes)}
\label{tab:failures}
\begin{tabular}{lccccc}
\toprule
\textbf{Model} & \textbf{Wrong Route} & \textbf{Wrong Urgency} & \textbf{Unhelpful Resp} & \textbf{No Follow-up} & \textbf{Step Limit} \\
\midrule
Claude 3.5 Sonnet & 0 & 0 & 1 & 1 & 0 \\
GPT-4o-Mini       & 1 & 1 & 2 & 2 & 0 \\
Gemini 2.0 Flash  & 1 & 2 & 3 & 3 & 0 \\
Llama-3.1-8B      & 6 & 4 & 7 & 5 & 0 \\
Mistral-7B        & 3 & 2 & 3 & 3 & 0 \\
\bottomrule
\end{tabular}
\end{table}

The primary bottleneck for smaller models is writing unhelpful responses (failing to include relevant troubleshooting steps) and failing to address follow-up customer emails. They also frequently mis-route tickets when incoming queries are ambiguous.

\subsection{Reward Hacking Analysis}

By logging occurrences where the keyword overlap score was high ($\ge 0.8$) but the LLM judge score was low ($< 0.4$), we measured the frequency of reward hacking. As shown in Figure~\ref{fig:hacking}, smaller models display much higher rates of reward hacking, often outputting key term lists rather than structured customer service letters. This suggests that smaller models struggle to balance formatting rules with conversational quality objectives.

\begin{figure}[ht]
\centering
\begin{tabular}{lc}
\toprule
\textbf{Model} & \textbf{Reward Hacking Rate (Flagged / Attempts)} \\
\midrule
Claude 3.5 Sonnet & 2.5\% (1/40) \\
GPT-4o-Mini       & 22.5\% (9/40) \\
Gemini 2.0 Flash  & 15.0\% (6/40) \\
Llama-3.1-8B      & 32.5\% (13/40) \\
Mistral-7B        & 42.5\% (17/40) \\
\bottomrule
\end{tabular}
\caption{Reward hacking rates detected during triage and resolution tasks.}
\label{fig:hacking}
\end{figure}

\subsection{Continuous Complexity Breakdown}

Mapping performance against our continuous ticket complexity buckets reveals that agent capabilities do not degrade discretely, but rather follow a smooth negative sigmoid curve. For instance, Claude 3.5 Sonnet maintains a high score of 0.947 on Low-complexity tickets ($C \le 0.4$) but falls to 0.739 on High-complexity tickets ($C > 0.7$). Mistral-7B scores 0.764 on Low complexity, but collapses to 0.400 on High complexity.

\section{Conclusion}

We introduced SupportOps, an OpenEnv benchmark representing a stateful customer triage workflow. By incorporating continuous difficulty scaling and a dual-signal grader, we show that SupportOps provides a highly calibrated environment for checking model capabilities and alignment issues like reward hacking. Future work will investigate using SupportOps to train agents via RL from AI Feedback (RLAIF) to actively suppress reward-hacking tendencies.

\end{document}
